\begin{table}[ht]
\centering
\small
\begin{tabular}{lcccc}
\toprule
 & & Human & \multicolumn{2}{c}{Projected saturation} \\
\cmidrule(lr){4-5}
Category & $n$ & baseline & Earliest & Latest \\
\midrule
Chemistry & 5 & 22--90\% & 2024-08 & 2026-01 \\
Cyber & 13 & --- & 2025-01 & 2027-05 \\
Trivia \& Commonsense QA & 5 & 80--96\% & 2022-10 & 2027-08 \\
Biology & 11 & 22--90\% & 2024-01 & 2027-12 \\
Autonomous SWE & 14 & 83\% & 2025-07 & 2028-01 \\
Core AGI Progress & 5 & 60--100\% & 2026-01 & 2028-02 \\
Advanced Language and Writing & 5 & --- & 2024-11 & 2028-02 \\
Domain Specific Questions & 8 & 22--94\% & 2023-06 & 2028-06 \\
Multimodal Understanding & 8 & 55--100\% & 2025-07 & 2028-08 \\
Mathematics & 8 & 35--97\% & 2023-09 & 2029-06 \\
Agentic Computer Use & 7 & 72\% & 2026-05 & 2029-07 \\
General Reasoning & 10 & 53--100\% & 2026-01 & 2030-01 \\
\bottomrule
\end{tabular}
\caption{\revised{\textbf{Projected saturation by capability category}, ordered by the latest date within the category. $n$ is the number of benchmarks in the category. The human-baseline column spans every human baseline recorded for the category's benchmarks, from crowdworkers to domain experts and expert committees, so a wide range reflects that mixture of populations rather than disagreement between measurements of the same quantity; Appendix~\ref{app:human_baselines} lists each baseline with its expertise level, and a dash means that no human baseline is available for any benchmark in the category. The last two columns give the earliest and latest posterior median saturation date among the category's benchmarks. Benchmarks are listed individually in Supp.\ Table~\ref{tab:benchmark_details} and described in Appendix~\ref{app:benchmarks}.}}
\label{tab:category_summary}
\end{table}
